\documentclass{article}

\usepackage{amsmath,amssymb}
\usepackage{xurl}
\usepackage[hidelinks]{hyperref}
\setlength{\emergencystretch}{2em}

% Shared commands for notes in docs/paper-gaps/.

\DeclareUnicodeCharacter{2080}{\ifmmode{}_{0}\else\textsubscript{0}\fi}
\DeclareUnicodeCharacter{2081}{\ifmmode{}_{1}\else\textsubscript{1}\fi}
\DeclareUnicodeCharacter{2082}{\ifmmode{}_{2}\else\textsubscript{2}\fi}
\DeclareUnicodeCharacter{2099}{\ifmmode{}_{n}\else\textsubscript{n}\fi}

\newcommand{\Lean}{\textsc{Lean}}
\ExplSyntaxOn
\NewDocumentCommand{\leanid}{m}{
  \ifmmode
    \textsf{\detokenize{#1}}
  \else
    \group_begin:
    \urlstyle{sf}
    \tl_set:Nn \l_tmpa_tl {#1}
    \tl_replace_all:Nnn \l_tmpa_tl {\_} {_}
    \exp_args:NV \path \l_tmpa_tl
    \group_end:
  \fi
}
\ExplSyntaxOff
\providecommand{\tr}{\operatorname{tr}}
\newcommand{\tnleanIssueBase}{https://github.com/LionSR/TNLean/issues/}
\newcommand{\tnleanPrBase}{https://github.com/LionSR/TNLean/pull/}
\newcommand{\tnleanDocBase}{https://sirui-lu.com/TNLean/docs/}
\newcommand{\ghissue}[1]{\href{\tnleanIssueBase#1}{GitHub issue \##1}}
\newcommand{\ghpr}[1]{\href{\tnleanPrBase#1}{PR \##1}}
\newcommand{\doclink}[1]{\href{\tnleanDocBase#1.html}{\path{#1}}}

% Dirac notation and the identity operator, as in blueprint/src/macros/common.tex.
% \providecommand keeps note-local definitions authoritative.
\usepackage{dsfont}
\providecommand{\ket}[1]{|#1\rangle}
\providecommand{\bra}[1]{\langle #1|}
\providecommand{\braket}[2]{\langle #1 | #2 \rangle}
\providecommand{\ketbra}[2]{|#1\rangle\!\langle #2|}
\providecommand{\Id}{\mathds{1}}
\providecommand{\one}{\mathds{1}}
\providecommand{\C}{\mathbb{C}}
\providecommand{\MN}[1]{M_{#1}(\C)}
\providecommand{\MD}{M_D(\C)}

% External referencing for paper sources.  The build helper writes sanitized
% auxiliary files under build/paper-aux/ and excludes duplicated source labels.
\usepackage{xr}

% Import a generated paper auxiliary file with a caller-supplied label prefix.
% The candidate paths support builds from docs/paper-gaps/, the repository root,
% and build/paper-aux/ itself.
\newcommand{\importpaperaux}[2]{%
  \IfFileExists{../../build/paper-aux/#2.aux}{%
    \externaldocument[#1]{../../build/paper-aux/#2}%
  }{%
    \IfFileExists{build/paper-aux/#2.aux}{%
      \externaldocument[#1]{build/paper-aux/#2}%
    }{%
      \IfFileExists{#2.aux}{%
        \externaldocument[#1]{#2}%
      }{}%
    }%
  }%
}

% General external reference macros
\newcommand{\paperref}[2]{\ref{#1-#2}}
\newcommand{\papereqref}[2]{(\ref{#1-#2})}

% CPSV16 source (1606.00608 / MPDO-22-12-17-2.tex) external document and shortcuts
\importpaperaux{cpsv16-}{1606.00608/MPDO-22-12-17-2-tnlean}
\newcommand{\cpsvref}[1]{\paperref{cpsv16}{#1}}
\newcommand{\cpsveqref}[1]{\papereqref{cpsv16}{#1}}

% Verdict marker: \gapnote{<kind>}{<status>}. Kind is the policy
% classification (clarification, local-correction, scope-restriction,
% unfaithful, false-source, open-gap); status is open, wip (elimination
% underway), resolved, or historical. Machine-read by the site index;
% prints a verdict line.
\newcommand{\gapnote}[2]{\par\noindent\textbf{Verdict:} #1 (#2).\par}


\title{The Square-Dimensional Scope of the Choi--Jamiolkowski Correspondence}
\author{}
\date{2026-07-30}

\begin{document}
\maketitle

\section{The source statement}

Let $T:\MN{d}\to\MN{d'}$ be linear and let
$\Omega_d\in\C^d\otimes\C^d$ be maximally entangled.  Wolf
Proposition~2.1 associates to $T$ the operator
\begin{align}
    \tau_T
      &= (T\otimes\operatorname{id}_d)
         \bigl(\ketbra{\Omega_d}{\Omega_d}\bigr)
       \in \MN{d'}\otimes\MN{d},
    \tag{Wolf 2.1}\\
    \tr\!\bigl(A\,T(B)\bigr)
      &= d\,\tr\!\bigl(\tau_T(A\otimes B^{T})\bigr)
    \qquad
    (A\in\MN{d'},\ B\in\MN{d}).
    \tag{Wolf 2.1}
\end{align}
These two formulas give mutually inverse correspondences between linear maps
and bipartite operators.  They identify Hermiticity preservation with
$\tau_T=\tau_T^\dagger$, complete positivity with $\tau_T\geq0$, unitality
with
\begin{align}
    \tr_B(\tau_T)&=\frac{\Id_{d'}}{d},
\end{align}
and trace preservation with
\begin{align}
    \tr_A(\tau_T)&=\frac{\Id_d}{d}.
\end{align}
More generally, without assuming trace preservation, the source gives
\begin{align}
    \tr_A(\tau_T)
      &= \frac{\bigl(T^*(\Id_{d'})\bigr)^T}{d},
    &
    \tr(\tau_T)
      &= \frac{\tr\!\bigl(T^*(\Id_{d'})\bigr)}{d}.
\end{align}
It also gives the corresponding proportional-identity criterion for a
doubly-stochastic map \cite[Proposition~2.1]{Wolf2012Quantum}.

\section{The square specialization}

The present Choi matrix takes a linear endomorphism
$T:\MN{D}\to\MN{D}$ and produces an operator on
$\C^D\otimes\C^D$.\footnote{Formal definition:
\leanid{ChoiJamiolkowski.choiMatrix} in
\doclink{TNLean/Channel/ChoiJamiolkowski}.}
Within this specialization, the development proves the complete-positivity
equivalence, the Hermiticity-preservation equivalence, and the implication
from trace preservation to
\begin{align}
    \tr_A(\tau_T)&=\frac{\Id_D}{D}.
\end{align}
It consequently proves $\tr(\tau_T)=1$ for trace-preserving maps.  It does not
yet prove the general adjoint partial-trace identity or the general trace
normalization identity displayed above.
Thus these results agree with Wolf Proposition~2.1 after imposing $d=d'$,
but they do not constitute its different-dimension statement.

\section{Missing clauses}

For arbitrary $d$ and $d'$, the remaining work is to construct both directions
of the correspondence, prove the trace-pairing formula, and derive the
Hermiticity, complete-positivity, unitality, trace-preservation, and
doubly-stochastic equivalences.  In particular, the converse partial-trace
implications, the general adjoint partial-trace identity, the general trace
normalization identity, and the proportional-identity clauses should be
established rather than inferred from the square specialization.

The natural complete-positivity predicate is the existence of rectangular
Kraus operators $A_i\in M_{d'\times d}(\C)$ satisfying
\begin{align}
    T(X)&=\sum_i A_i X A_i^\dagger.
\end{align}
Using this predicate keeps the Choi correspondence compatible with the
rectangular Kraus representation already used elsewhere in the development.

\section{The related Kraus-representation restriction}

Wolf Theorem~2.1 states that a map
$T:\MN{d}\to\MN{d'}$ is completely positive if and only if it has a
rectangular Kraus representation
\begin{align}
    T(X)&=\sum_{j=1}^{r}K_jXK_j^\dagger,
    &K_j&\in M_{d'\times d}(\C).
    \tag{Wolf 2.8}
\end{align}
It further identifies the trace-preserving and unital normalizations, the
minimal value $r=\operatorname{rank}(\tau_T)\leq dd'$, an orthogonal minimal
family, and the unitary freedom after padding the smaller family with zero
operators \cite[Theorem~2.1]{Wolf2012Quantum}.

The present development contains rectangular Kraus predicates and several
normalization and freedom results, but not this full arbitrary-dimensional
package.  The missing work is the rectangular existence equivalence together
with the minimal-rank and orthogonal-family conclusions, followed by the
padded-unitary classification in the source's stated generality.

\section{Verdict}

These are \emph{scope restrictions}: the proved Choi and Kraus results are
mathematically correct in their stated settings, while Wolf Proposition~2.1
and Theorem~2.1 allow independent input and output dimensions and include the
additional clauses above.  The source-labelled blueprint entries therefore
state the restrictions explicitly.  The restrictions can be removed once the
different-dimension statements have been proved.\footnote{Formal follow-up:
\ghissue{3987}.}

\bibliographystyle{alpha}
\bibliography{references}

\end{document}
